\notechapter{N-20221222}{Topological Re-definitions, Subspaces, Connectedness, and Path-Connectedness}

\section{Re-defining continuity}

After defining topological spaces, the note rewrites metric definitions using open sets.  A function \(f:X\to Y\) is continuous if
\[
  f^{-1}(V)\text{ is open in }X
\]
for every open \(V\subseteq Y\).  This definition no longer mentions \(\varepsilon\), \(\delta\), or distances.  It says that open information in \(Y\) pulls back to open information in \(X\).

A homeomorphism is a bijection preserving open sets in both directions.  This is the topological notion of sameness.

\section{Closed sets and closure}

A set \(C\subseteq X\) is closed if \(X\setminus C\) is open.  The closure of \(S\) is the smallest closed set containing \(S\):
\[
  \overline S=\bigcap\{C:C\text{ closed and }S\subseteq C\}.
\]
The note warns that complete and totally bounded are metric properties, not topological properties.  It also warns that sequences do not work perfectly in general topology; one often needs nets or filters.

\section{Subspaces}

If \(Y\subseteq X\), the subspace topology on \(Y\) is
\[
  \tau_Y=\{U\cap Y:U\in\tau_X\}.
\]
The key warning is: open in \(Y\) need not be open in \(X\).  For example,
\[
  [0,1)=[0,2]\cap(-1,1)
\]
is open in \([0,2]\), but not open in \(\mathbb R\).

\section{Connectedness}

A space \(X\) is disconnected if there are nonempty open sets \(U,V\) such that
\[
  X=U\cup V,\qquad U\cap V=\varnothing.
\]
If no such separation exists, \(X\) is connected.  Equivalently, \(X\) has no nontrivial clopen subset.  Separated unions such as \([1,2]\cup[3,4]\) are disconnected, while intervals are connected.

\section{Path-connectedness}

A space \(X\) is path-connected if for any \(x,y\in X\), there is a continuous map
\[
  \gamma:[0,1]\to X,\qquad \gamma(0)=x,\quad \gamma(1)=y.
\]
Path-connectedness implies connectedness.  The converse can fail.  In ordinary regions of \(\mathbb R^n\), it matches the intuition that one can walk from one point to another without leaving the set.


\section{Continuity: why preimages appear}

The preimage definition of continuity can feel backward at first.  We define continuity by saying that for every open set \(V\subseteq Y\), the preimage \(f^{-1}(V)\) is open in \(X\), not by saying images of open sets are open.  The reason is logical: to know whether \(f(x)\) lies in a neighborhood \(V\), we ask which points \(x\) are sent into \(V\).  That set is the preimage.

Images of open sets need not be open.  For example, the map
\[
  f:\mathbb R\to\mathbb R,\qquad f(x)=x^2
\]
sends the open interval \((-1,1)\) to \([0,1)\), which is not open in \(\mathbb R\).  But \(f\) is continuous, because preimages of open sets are open.

This example is important because it prevents a common misconception: continuity is about pulling back open tests, not pushing open sets forward.

\section{Closure through neighborhoods}

The closure \(\overline S\) can also be described pointwise:
\[
  x\in\overline S
  \quad\Longleftrightarrow\quad
  \text{every open neighborhood of }x\text{ meets }S.
\]
This description is often easier than the intersection-of-closed-sets definition.  For instance,
\[
  \overline{\mathbb Q}=\mathbb R
\]
in the usual topology because every interval contains rational numbers.  Similarly,
\[
  \overline{\mathbb R\setminus\mathbb Q}=\mathbb R.
\]
The two dense subsets are disjoint, showing that closure is not about having large measure or containing intervals; it is about meeting every neighborhood.

\section{Connectedness by continuous images}

A useful theorem is that continuous images of connected spaces are connected.  If \(X\) is connected and \(f:X\to Y\) is continuous, then \(f(X)\) is connected.  Otherwise, if \(f(X)=U\cup V\) were a separation, then \(f^{-1}(U)\) and \(f^{-1}(V)\) would separate \(X\).

This theorem proves that intervals in \(\mathbb R\) are connected.  If an interval were split into two nonempty open pieces, one could construct a continuous map to the discrete two-point space, contradicting the intermediate value principle.  Conversely, the intermediate value theorem is essentially a connectedness statement about intervals.

\section{Path-connectedness and components}

Path-connectedness is stronger because it requires an explicit continuous path.  The path-connected component of a point \(x\) is the set of all points reachable from \(x\) by paths.  These components partition the space.

Every path-connected space is connected.  The proof is simple: if \(X=U\cup V\) were a separation and a path \(\gamma:[0,1]\to X\) started in \(U\) and ended in \(V\), then \([0,1]\) would be separated by \(\gamma^{-1}(U)\) and \(\gamma^{-1}(V)\), impossible because \([0,1]\) is connected.

The converse fails in spaces like the topologist's sine curve.  That example is not needed for computation, but it explains why topology keeps the two definitions separate.

\section{Next viewpoint}

Continuity, closure, subspace topology, connectedness, and path-connectedness are open-set notions.  Boundedness and completeness are not.  The note is teaching what topology remembers and what it deliberately forgets.

\section{Continuity as a contravariant test}

The preimage definition of continuity has a categorical flavor: open sets pull back.  If
\[
  X\xrightarrow{f}Y\xrightarrow{g}Z
\]
are continuous, then for every open \(W\subseteq Z\),
\[
  (g\circ f)^{-1}(W)=f^{-1}(g^{-1}(W))
\]
is open in \(X\).  This one-line identity is why continuous maps compose.

Images do not behave this cleanly.  A continuous map may send open sets to non-open sets and closed sets to non-closed sets.  Preimages are the stable operation; topology is built around that stability.

\section{Closure as an operator}

The closure operation satisfies Kuratowski-style rules:
\[
  S\subseteq \overline S,\qquad
  \overline{\overline S}=\overline S,\qquad
  \overline{S\cup T}=\overline S\cup\overline T.
\]
It also preserves inclusions: if \(S\subseteq T\), then \(\overline S\subseteq\overline T\).

These algebraic properties turn closure into more than a point-set definition.  They let one reason with closure abstractly, without drawing neighborhoods each time.

\section{Why sequences are not enough}

In metric spaces, closure can be detected by sequences.  In arbitrary topological spaces, this may fail.  The replacement notions are nets and filters.  A net is like a sequence whose index set is not necessarily \(\mathbb N\), but any directed set.  This extra flexibility lets nets test all neighborhoods in spaces that are too large or too non-countable for ordinary sequences.

The warning is not meant to make topology harder for its own sake.  It explains exactly which metric-space proof habits are safe and which are not.  Sequential reasoning works in first-countable spaces, including metric spaces, but general topology requires a broader convergence language.

\section{Connectedness as maps to a two-point space}

A space \(X\) is disconnected iff there is a nonconstant continuous map
\[
  X\to \{0,1\}
\]
where \(\{0,1\}\) has the discrete topology.  Indeed, the preimages of \(0\) and \(1\) are disjoint nonempty clopen sets.

So connectedness can be read as: every continuous discrete-valued function is constant.  This perspective is useful because it turns a separation problem into a mapping problem.

\section{Path components versus connected components}

Path components always lie inside connected components.  Equality holds in many familiar spaces, such as open subsets of \(\mathbb R^n\), because they are locally path-connected.  But equality can fail in pathological spaces such as the topologist's sine curve.

The difference matters conceptually.  Connectedness says the space cannot be split by open sets.  Path-connectedness says points can be joined by a continuous interval.  The first is global and negative; the second is constructive.

\section{Local connectedness as the missing hypothesis}

A space is locally path-connected if every point has a basis of path-connected neighborhoods.  In such spaces, connected components and path components agree, and path components are open.  This is why ordinary geometric regions behave better than arbitrary topological spaces.

Thus the hierarchy is:
\[
  \text{path-connected}\Rightarrow\text{connected},
\]
but local path-connectedness often lets one recover path information from connectedness.

\section{Continuity, closure, and connectedness}

This note is about what survives after distances disappear.  Continuity is preimage stability, closure is unavoidable nearness, subspace topology is restriction, connectedness is the impossibility of a clopen split, and path-connectedness is reachability by intervals.  The advanced warning is that metric intuition remains useful but not universal: sequences may need nets, and connectedness may not mean path-connectedness.
